\documentclass[11pt,a4paper]{article}
\usepackage[utf8]{inputenc}
\usepackage[T1]{fontenc}
\usepackage[spanish,es-nodecimaldot]{babel}
\usepackage{lmodern}
\usepackage{amsmath}
\usepackage{booktabs}
\usepackage{geometry}
\usepackage{xcolor}
\usepackage{hyperref}
\usepackage{listings}
\usepackage{microtype}
\geometry{margin=2.35cm}
\hypersetup{colorlinks=true,linkcolor=blue!55!black,urlcolor=blue!55!black}
\lstdefinestyle{python}{language=Python,basicstyle=\ttfamily\small,
  keywordstyle=\color{blue!65!black},commentstyle=\color{green!45!black},
  stringstyle=\color{orange!55!black},breaklines=true,frame=single,
  showstringspaces=false,columns=fullflexible}
\lstdefinestyle{yaml}{basicstyle=\ttfamily\small,breaklines=true,frame=single,
  showstringspaces=false,columns=fullflexible}

\title{Explicación del código fuente de locomoción\\Nova Spot Micro}
\author{Proyecto de tesis --- robot cuadrúpedo Nova Spot Micro}
\date{Actualizado el 16 de agosto de 2026}

\begin{document}
\maketitle
\begin{abstract}
Este documento explica cómo el código ROS 2 genera la marcha cartesiana, aplica
cinemática inversa y controla las doce articulaciones del Nova Spot Micro en
Gazebo y MuJoCo. Se describen responsabilidades, flujo de datos, validaciones,
configuración, pruebas y comandos de ejecución. El controlador es todavía una
línea base de simulación; no debe conectarse directamente a servomotores sin
calibración y supervisor físico.
\end{abstract}
\tableofcontents

\section{Arquitectura general}
El sistema se divide en dos paquetes ROS 2:
\begin{itemize}
\item \texttt{nova\_sm3\_description}: URDF/Xacro, propiedades físicas
  provisionales, integración \texttt{ros2\_control} y lanzador de Gazebo.
\item \texttt{nova\_gait\_controller}: cinemática, generación de marchas,
  nodo de control, parámetros y pruebas.
\end{itemize}

El flujo de ejecución es
\[
\text{comando ``gateo''}\rightarrow
\text{trayectoria de pies}\rightarrow
\text{IK}\rightarrow
\texttt{JointTrajectory}\rightarrow
\texttt{ros2\_control}\rightarrow\text{Gazebo}.
\]
Gazebo devuelve los estados articulares mediante \texttt{/joint\_states}; la
pose del cuerpo, la fase y los contactos medidos se publican para validar
altura, orientación y cumplimiento del patrón de apoyo.

\section{Mapa de archivos}
\begin{table}[h]
\centering
\small
\begin{tabular}{p{0.42\textwidth}p{0.50\textwidth}}
\toprule
Archivo & Responsabilidad\\
\midrule
\texttt{nova\_gait\_controller/kinematics.py} & FK, IK, límites y generación
cartesiana del gateo y de la marcha paso.\\
\texttt{nova\_gait\_controller/gait\_controller.py} & Nodo ROS 2, máquina de
estados, parámetros, fase y publicación de trayectorias.\\
\texttt{contact\_monitor.py}, \texttt{contact\_comparator.py} & Agregación de
sensores y comparación de contactos previstos frente a observados.\\
\texttt{metrics\_node.py}, \texttt{safety\_supervisor.py} & Métricas de pose y
articulaciones; parada lógica por altura o inclinación insegura.\\
\texttt{config/gaits.yaml} & Valores operativos de marcha.\\
\texttt{test/test\_kinematics.py} & Pruebas FK--IK y continuidad.\\
\texttt{nova\_sm3\_description/urdf/*.xacro} & Geometría, articulaciones,
límites e interfaces de control.\\
\texttt{nova\_sm3\_description/config/controllers.yaml} & Configuración del
administrador y controladores.\\
\texttt{nova\_sm3\_description/launch/sim.launch.py} & Inicio de Gazebo,
publicación del modelo y activación de controladores.\\
\texttt{nova\_gait\_controller/launch/demo.launch.py} & Composición de
simulador y nodo de marcha.\\
\bottomrule
\end{tabular}
\end{table}

\section{Módulo \texttt{kinematics.py}}
\subsection{Constantes geométricas}
El archivo declara las tres longitudes que deben coincidir con el URDF:
\begin{lstlisting}[style=python]
COXA_LENGTH = 0.090
FEMUR_LENGTH = 0.105
TIBIA_LENGTH = 0.132
\end{lstlisting}
Se mantienen centralizadas para evitar mezclar las dimensiones del repositorio
de referencia con las del NovaSM3. En una fase posterior deben leerse de una
fuente única generada a partir de mediciones.

\subsection{\texttt{forward\_leg}}
Recibe tres ángulos y \texttt{side}, que vale $+1$ a la izquierda y $-1$ a la
derecha. Primero calcula la cadena fémur--tibia en el plano y después aplica la
rotación de la coxa al plano lateral/vertical. Devuelve una tupla
\texttt{(x, y, z)} relativa al pivote de cadera.

Esta función se utiliza para obtener la ubicación neutral de cada pie desde la
postura nominal. También permite comprobar que la IK reconstruye los ángulos
de partida.

\subsection{\texttt{inverse\_leg}}
Recibe una posición de pie y devuelve la terna articular compatible con los
ejes del URDF. El orden de sus defensas es importante:
\begin{enumerate}
\item verifica que exista separación radial suficiente para la coxa;
\item calcula y normaliza el ángulo de coxa;
\item aplica ley de cosenos al plano fémur--tibia;
\item limita errores numéricos pequeños de $D$ al intervalo $[-1,1]$;
\item selecciona la rama de rodilla negativa;
\item rechaza valores no finitos o fuera de límites.
\end{enumerate}
Los errores se comunican mediante \texttt{ValueError}. El controlador captura
esta excepción antes de activar la marcha.

\subsection{\texttt{cartesian\_crawl}}
Esta función genera toda la tabla articular de un ciclo. Sus argumentos son la
postura, el número de muestras, la longitud y la elevación del paso. Los
desfases de pata son:
\begin{lstlisting}[style=python]
offsets = {'fl': 0.75, 'rr': 0.50,
           'fr': 0.25, 'rl': 0.00}
\end{lstlisting}
Al tener cuatro desfases separados por un cuarto de ciclo y un apoyo del 75\%,
solamente una pata se eleva. En apoyo, el pie retrocede linealmente respecto al
cuerpo; en oscilación, avanza y sigue una parábola vertical. Cada punto se
convierte inmediatamente con \texttt{inverse\_leg}. El resultado es una lista
de $N$ vectores, cada uno con doce ángulos en el orden exigido por ROS 2.

La medición de contacto reveló que esta formulación no realiza transferencia
lateral explícita: aunque el cuerpo avanza, las patas traseras permanecen
apoyadas y deslizan durante su oscilación prevista. Esta limitación motiva una
nueva versión de la trayectoria.

\subsection{\texttt{cartesian\_step\_walk}}
Genera 32 referencias por ciclo con paso de 16 mm, elevación de 8 mm y
transferencia lateral sinusoidal de 4 mm. Mantiene el orden FL--RR--FR--RL y
fue validada en dos ensayos de Gazebo y una ejecución headless de MuJoCo.

\section{Nodo \texttt{gait\_controller.py}}
\subsection{Orden articular}
La constante \texttt{JOINTS} fija el contrato con
\texttt{joint\_trajectory\_controller}: delantera izquierda, delantera derecha,
trasera izquierda y trasera derecha; en cada pata, coxa, fémur y tibia. Alterar
este orden sin cambiar el URDF y la configuración movería articulaciones
equivocadas.

\subsection{Postura nominal}
\begin{lstlisting}[style=python]
STAND = (0.10, 0.42, -0.84)
\end{lstlisting}
La función \texttt{pose} repite o combina cuatro ternas y las aplana en un
vector de doce elementos. \texttt{stand} y \texttt{stop} publican esta postura
en una transición de $0.8$ s. En la implementación actual \texttt{stop} detiene
la secuencia y conserva postura; no equivale a cortar físicamente la energía.

\subsection{Inicialización ROS 2}
El constructor:
\begin{enumerate}
\item declara parámetros;
\item crea un publicador de \texttt{JointTrajectory};
\item se suscribe a \texttt{/nova/gait\_command};
\item inicializa modo, fase y reloj;
\item crea un temporizador de 20 ms;
\item entra en la marcha inicial, normalmente \texttt{stand}.
\end{enumerate}

\subsection{Comandos}
\texttt{command\_callback} normaliza alias en español:
\texttt{gateo}$\rightarrow$\texttt{crawl},
\texttt{paso}$\rightarrow$\texttt{step},
\texttt{galope}$\rightarrow$\texttt{gallop} y
\texttt{parar}$\rightarrow$\texttt{stop}. \texttt{set\_mode} rechaza nombres
desconocidos y reinicia el índice de fase cuando acepta un cambio.
El galope se rechaza mientras
\texttt{enable\_experimental\_gallop=false}; solo se permite una habilitación
explícita para experimentos opcionales de simulación y no para hardware.

Al solicitar \texttt{crawl}, \texttt{build\_crawl} lee los parámetros actuales,
valida sus intervalos y construye las referencias cartesianas. Esto permite
cambiar el YAML o parámetros ROS sin modificar las ecuaciones. Si la tabla no
puede construirse, se registra el error, se conserva \texttt{stand} y no se
publica el gateo.

\subsection{Bucle temporal}
El temporizador llama a \texttt{update}. La función retorna sin publicar si el
modo está detenido o todavía no se alcanza el instante de la fase siguiente.
Cuando corresponde, selecciona la tabla, publica el vector y un JSON sincronizado
en \texttt{/nova/gait\_phase}, incrementa el índice y programa el siguiente
instante desde el vencimiento planificado anterior, sin acumular jitter.

\subsection{Publicación de trayectorias}
\texttt{publish\_target} crea un único \texttt{JointTrajectoryPoint}:
\begin{itemize}
\item copia las doce posiciones para no compartir referencias mutables;
\item solicita velocidad final cero en las doce articulaciones;
\item convierte el tiempo decimal en segundos y nanosegundos ROS;
\item adjunta los nombres articulares y publica el mensaje.
\end{itemize}
El parámetro \texttt{transition\_ratio=0.90} hace que la interpolación termine
antes del límite completo de fase.

\section{Configuración \texttt{gaits.yaml}}
La línea base validada es:
\begin{lstlisting}[style=yaml]
gait_controller:
  ros__parameters:
    initial_gait: stand
    crawl_phase_duration: 0.18
    crawl_samples: 24
    crawl_step_length: 0.018
    crawl_step_height: 0.014
    transition_ratio: 0.90
\end{lstlisting}
Los intervalos admitidos están codificados en \texttt{build\_crawl}. Que un
valor esté dentro del intervalo no garantiza por sí solo estabilidad en todas
las combinaciones; la IK y las pruebas de simulación son defensas adicionales.

\section{Descripción, controladores y Gazebo}
El Xacro declara doce articulaciones controladas con interfaz de comando de
posición y estados de posición/velocidad. Cada articulación recibe como valor
inicial la postura nominal, evitando una caída libre antes de activar el
controlador.

\texttt{controllers.yaml} configura:
\begin{itemize}
\item \texttt{joint\_state\_broadcaster}, que publica estados;
\item \texttt{joint\_trajectory\_controller}, que interpola y aplica los puntos.
\end{itemize}

\texttt{sim.launch.py} procesa el Xacro con Gazebo habilitado, inicia el mundo
vacío, publica \texttt{robot\_description}, crea el robot a $z=0.245$ m y,
después de cinco segundos, activa ambos controladores. El lanzador
\texttt{demo.launch.py} añade el nodo de marcha usando el YAML.

\section{Pruebas automatizadas}
La suite cubre FK--IK, continuidad y periodicidad de las marchas, temporización,
plan de contactos, diagnósticos, seguridad, dinámica y consistencia
Python--URDF--MJCF. La verificación actual registra 33 pruebas aprobadas. Además, los
módulos pasan \texttt{py\_compile}, el URDF pasa \texttt{check\_urdf} y ambos
paquetes compilan con \texttt{colcon}.

\section{Ejecución reproducible}
\begin{lstlisting}[style=yaml]
cd /home/pavilion/Documentos/Cuadrupedo
source /opt/ros/humble/setup.bash
colcon build --symlink-install
source install/setup.bash
ros2 launch nova_gait_controller demo.launch.py
\end{lstlisting}
En otra terminal:
\begin{lstlisting}[style=yaml]
source /opt/ros/humble/setup.bash
source /home/pavilion/Documentos/Cuadrupedo/install/setup.bash
ros2 topic pub --once /nova/gait_command \
  std_msgs/msg/String "{data: gateo}"
\end{lstlisting}
Comandos de postura y detención lógica:
\begin{lstlisting}[style=yaml]
ros2 topic pub --once /nova/gait_command \
  std_msgs/msg/String "{data: stand}"
ros2 topic pub --once /nova/gait_command \
  std_msgs/msg/String "{data: stop}"
\end{lstlisting}

\section{Resultado y limitaciones}
La cadencia corregida del gateo se reprodujo en dos ensayos de 13 ciclos con
periodo medio cercano a 4.320 s y diferencias menores al 0.2\% en las medias.
La marcha paso completó dos ensayos de 12 ciclos en Gazebo y 12 ciclos en
MuJoCo. El sistema publica métricas, posee un supervisor provisional y mide los
cuatro contactos en Gazebo.

El ensayo cuantitativo de 76 ciclos obtuvo solo 32.550\% de coincidencia de
contactos: las patas delanteras tuvieron transiciones tardías y las traseras no
despegaron. Faltan IMU simulada, contactos equivalentes en MuJoCo, margen de
estabilidad en línea e interfaz física PCA9685.
Las constantes geométricas y dinámicas son provisionales. Por ello, este código
es válido como línea base de simulación y no autoriza energizar los MG996R.

\section{Próximas extensiones recomendadas}
\begin{enumerate}
\item corregir el gateo y repetir el ensayo cuantitativo de contactos;
\item añadir contactos equivalentes en MuJoCo y una IMU simulada;
\item calcular polígono de soporte y margen del centro de masa;
\item ampliar y rearmar de forma segura el supervisor;
\item sustituir longitudes, límites y dinámica por mediciones físicas;
\item mantener esta marcha como referencia nominal para comparar posteriormente
  la capa correctiva de aprendizaje por refuerzo.
\end{enumerate}

\end{document}
